\chapter{Ảnh vệ tinh}
\label{ch:M6L2LessonNotes}

\section*{Lesson Notes}

MÔ-ĐUN 6 \textbar{} BÀI 2

\section{Tài liệu đọc bắt buộc}

Các tài liệu đọc bắt buộc của chương trình này đều là tài liệu truy cập mở, nghĩa là bạn có thể truy cập miễn phí. Liên kết trong phần trích dẫn sẽ đưa bạn trực tiếp đến tài liệu hoặc đến trang cho phép tải xuống.

\begin{enumerate}
\def\labelenumi{\arabic{enumi}.}
\item
  NASA. "Remote Sensing." \emph{NASA EarthData}. https://www.earthdata.nasa.gov/learn/backgrounders/remote-sensing\#spectral-resolution
\end{enumerate}

\begin{itemize}
\item
  \textbf{Các trang bắt buộc:} Toàn bộ trang web
\item
  \textbf{Thời gian đọc ước tính:} 30 phút
\end{itemize}

\begin{enumerate}
\def\labelenumi{\arabic{enumi}.}
\setcounter{enumi}{1}
\item
  NASA Goddard. "9 Things about Landsat 9." \emph{YouTube}, đăng tải ngày 26 tháng 9 năm 2021, https://www.youtube.com/watch?v=DGE-N8\_LQBo.
\end{enumerate}

\begin{itemize}
\item
  \textbf{Các trang bắt buộc:} Toàn bộ video
\item
  \textbf{Thời gian xem ước tính:} 6 phút
\end{itemize}

MODULE 6 \textbar{} BÀI 2

\section{Ảnh vệ tinh}

{\small
\begin{longtable}[]{@{}ll@{}}
\toprule\noalign{}
\endhead
\bottomrule\noalign{}
\endlastfoot
\textbf{Thời gian đọc} & 60 phút \\
\textbf{Kiến thức nền} & Python cơ bản \\
\textbf{Từ khóa} & Viễn thám, Phổ điện từ, Độ phân giải không gian \\
& \\
& \\
\end{longtable}
}

\section{Giới thiệu}

Trong bài học này, chúng ta sẽ làm quen với một loại dữ liệu mà chúng ta chưa đề cập: \textbf{ảnh vệ tinh}. Theo truyền thống, ảnh vệ tinh là loại dữ liệu được sử dụng trong các ngành khoa học tự nhiên hoặc lĩnh vực quốc phòng. Tuy nhiên, trong những năm gần đây, ngày càng nhiều nhà phân tích tài chính sử dụng ảnh vệ tinh để phân tích, với hy vọng thu được thêm những hiểu biết bổ sung bên cạnh phân tích tài chính truyền thống để đưa ra các quyết định tài chính tốt hơn. Nội dung của bài học này dựa trên thông tin trong các tài liệu đọc bắt buộc của bạn. Trước hết, chúng ta sẽ giới thiệu những kiến thức cơ bản về ảnh vệ tinh. Sau đó, chúng ta sẽ trình bày ngắn gọn một số ứng dụng của ảnh vệ tinh trong phân tích tài chính. Cuối cùng, chúng ta sẽ sử dụng Google Earth Engine (GEE) để truy cập một số ảnh vệ tinh và đưa ra các ví dụ về cách sử dụng ảnh vệ tinh cho phân tích.

\section{Ứng dụng tài chính của ảnh vệ tinh}

Trong phần này, chúng ta sẽ tìm hiểu cách đưa ảnh vệ tinh vào làm một phần dữ liệu cho phân tích tài chính.

Một trong những trường hợp nổi tiếng nhất về việc sử dụng ảnh vệ tinh trong phân tích tài chính là đếm số lượng ô tô ra vào bãi đỗ xe của các cửa hàng Walmart trong một khoảng thời gian. Các nhà phân tích tài chính dùng thông tin này như một chỉ báo thay thế (proxy) cho lưu lượng khách đến cửa hàng để ước tính doanh số của cửa hàng trong một giai đoạn cụ thể. Đây chính là ứng dụng đầu tiên của ảnh vệ tinh.

\subsection{Giám sát hoạt động kinh tế của một địa điểm}

Ví dụ về bãi đỗ xe của Walmart là một ứng dụng kinh điển của ảnh vệ tinh. Một ví dụ khác là dùng ảnh vệ tinh để giám sát lưu lượng tàu container tại một kênh đào lớn hoặc một cảng lớn nhằm phát hiện các tắc nghẽn có thể xảy ra trong chuỗi cung ứng. Chúng ta cũng có thể dùng ảnh vệ tinh để giám sát các nhà máy nhằm đánh giá mức sản lượng tiềm năng tại một thời điểm nhất định. Ngoài ra, chúng ta còn có thể giám sát một mỏ khai thác để theo dõi sản lượng hàng hóa cơ bản (commodity).

\subsection{Theo dõi sản lượng cây trồng}

Chúng ta có thể dùng ảnh vệ tinh để xem một loại cây trồng nhất định, như lúa mì hoặc ngô, có khỏe mạnh hay không. Từ thông tin này, các nhà phân tích tài chính có thể ước tính thị trường sẽ dư cung hay thiếu hụt nguồn cung của loại cây trồng đó và giao dịch dựa trên thông tin này. Thông thường, chúng ta sẽ dùng kết hợp các dải hồng ngoại để thực hiện phân tích này.

\subsection{Điều tra thiệt hại do thiên tai}

Hiện nay, các công ty bảo hiểm sử dụng ảnh vệ tinh để đánh giá thiệt hại do lũ lụt hoặc cháy rừng. Phương pháp thu thập thông tin thiệt hại này hiệu quả hơn so với phương pháp truyền thống là đến tận hiện trường. Thứ nhất, phương pháp truyền thống có thể cần vài tháng để thu thập dữ liệu, trong khi dùng ảnh vệ tinh chỉ mất vài giờ. Thứ hai, nhiều khu vực bị thiệt hại thường không thể tiếp cận được sau thảm họa, nên không thể áp dụng phương pháp truyền thống. Với công nghệ ảnh vệ tinh mới này, các công ty bảo hiểm có thể xử lý yêu cầu bồi thường nhanh hơn và tiết kiệm chi phí hơn.

\subsection{Những hạn chế khi sử dụng ảnh vệ tinh cho phân tích tài chính}

Hiện nay, đối với công chúng nói chung, có một số vấn đề khi sử dụng ảnh vệ tinh cho phân tích đầu tư.

~~~(A) Độ phân giải thời gian của ảnh vệ tinh công khai có thể dài. Ví dụ, Landsat 9 có độ phân giải thời gian là 8 ngày. Nhiều vệ tinh thương mại cung cấp ảnh vệ tinh với độ phân giải thời gian ngắn, chẳng hạn một ngày. Tuy nhiên, chúng rất đắt đỏ.

~~~(B) Xử lý lượng lớn ảnh vệ tinh đòi hỏi hạ tầng máy tính mạnh, thứ mà cá nhân thường không có.

~~~(C) Xử lý ảnh vệ tinh thô và chuyển chúng thành ảnh có thể sử dụng đòi hỏi được đào tạo chuyên môn.

Do những hạn chế này khi dùng ảnh vệ tinh để phân tích, hiện có nhiều công ty cung cấp dịch vụ truy cập các nền tảng phân tích ảnh vệ tinh của họ theo hình thức đăng ký thuê bao. Các công ty này thu thập ảnh vệ tinh công khai và tư nhân, xử lý chúng, rồi tải lên nền tảng của mình. Các nền tảng này thường tích hợp khả năng học máy để phân tích. Tuy nhiên, phí truy cập các nền tảng này cao, nên chỉ các quỹ phòng hộ (hedge fund) hoặc các công ty đầu tư lớn mới có thể sử dụng. Lấy ví dụ việc dùng lưu lượng ô tô tại bãi đỗ xe của các cửa hàng bán lẻ làm chỉ báo cho giao dịch cổ phiếu. Đây vẫn là một chiến lược giao dịch sinh lời, nhưng chỉ các quỹ phòng hộ mới tiếp cận được thông tin này.

\section{Kiến thức cơ bản về ảnh vệ tinh}

\subsection{Viễn thám}

Trước khi bàn về ảnh vệ tinh, chúng ta cần giới thiệu một vài khái niệm. Khái niệm đầu tiên là viễn thám. \textbf{Viễn thám} (remote sensing) là hoạt động sử dụng các vệ tinh quan sát Trái Đất hoặc máy bay được trang bị cảm biến để thu nhận hình ảnh và các thông tin khác về bề mặt Trái Đất từ trên cao.

Các cảm biến trên vệ tinh đo \textbf{bức xạ điện từ (EMR)}, vốn trước hết được phát ra từ Mặt Trời. Khi đến bề mặt Trái Đất, bức xạ này được phản xạ trở lại các cảm biến trên vệ tinh. Hình 1 dưới đây minh họa đường đi của EMR.

\textbf{Hình 1: Minh họa cách vệ tinh đo EMR phản xạ từ bề mặt Trái Đất}

Phỏng theo: NASA EarthData. "What Is Remote Sensing?" 10 December 2024. https://www.earthdata.nasa.gov/learn/backgrounders/remote-sensing.

\subsection{Phổ điện từ}

EMR là một dạng năng lượng truyền đi theo sóng. Tuy nhiên, các sóng này không đồng nhất; chúng có các bước sóng khác nhau. Bước sóng của một sóng được đo bằng khoảng cách giữa hai đỉnh sóng. Nếu một sóng có bước sóng ngắn hơn các sóng khác thì sóng này có tần số cao hơn, vì sóng dao động lên xuống rồi lại lên nhanh hơn.

Chúng ta nhóm các sóng có bước sóng tương tự nhau thành một dải gọi là \textbf{phổ (spectrum)}. Hình 2 bên dưới cho thấy cách chúng ta nhóm các sóng thành các phổ hay dải khác nhau theo bước sóng của chúng.

\textbf{Hình 2: Phổ bức xạ điện từ}

Nguồn: NASA EarthData. "What Is Remote Sensing?" 10 December 2024. https://www.earthdata.nasa.gov/learn/backgrounders/remote-sensing.

Hình 2 biểu diễn toàn bộ phổ bức xạ điện từ. Tuy nhiên, nhìn chung, các vệ tinh quan sát Trái Đất không thu thập dữ liệu trên toàn bộ phổ. Vệ tinh thường chỉ thu thập thông tin từ các dải hồng ngoại và các dải mà mắt người nhìn thấy được. Thông tin từ các dải này cung cấp đủ dữ liệu để chúng ta khảo sát các hoạt động diễn ra trên bề mặt Trái Đất.

Mọi bức xạ điện từ đều là ánh sáng, nhưng chỉ một số bước sóng nhất định mới được mắt người phát hiện. Đây được gọi là phổ ánh sáng khả kiến. Mọi màu sắc mà con người nhìn thấy đều là tổ hợp của ba màu cơ bản: đỏ, lục và lam (RGB). Mỗi màu cơ bản này đại diện cho một dải.

Ảnh chụp từ máy ảnh cũng là tổ hợp của ba màu cơ bản này. Các tổ hợp này có thể cho thấy các mức thông tin khác nhau về Trái Đất, và rất hữu ích đối với các nhà nghiên cứu. Chúng ta sẽ chỉ sử dụng một vài tổ hợp trong bài học này. Với những ai muốn tìm hiểu thêm về các tổ hợp, có rất nhiều thông tin trên internet.

Các vệ tinh thu thập thông tin từ nhiều dải được gọi là vệ tinh đa phổ (multispectral). Mỗi vệ tinh có thể xây dựng hệ thống dải riêng của mình. Ví dụ, một vệ tinh phổ biến mà công chúng có thể truy cập ảnh vệ tinh miễn phí là các vệ tinh Landsat của Hoa Kỳ. Hình 3 bên dưới cho thấy cấu trúc dải của vệ tinh Landsat 8 của Hoa Kỳ.

\textbf{Hình 3: Cấu trúc dải của vệ tinh Landsat 8 của Hoa Kỳ}

Nguồn: USGS. "What Are the Band Designations for the Landsat Satellites?" 31 May 2024. https://www.usgs.gov/faqs/what-are-band-designations-landsat-satellites

Theo Hình 3, nếu muốn tạo ra một ảnh vệ tinh trông giống ảnh chụp từ máy ảnh, chúng ta cần trích xuất ảnh từ dải 4 (đỏ), dải 3 (lục) và dải 2 (lam) rồi kết hợp chúng lại thành một ảnh giống ảnh chụp thông thường. Hình 4 bên dưới minh họa khái niệm này.

\textbf{Hình 4: Các dải màu để tạo ảnh giống ảnh chụp}

Nguồn: Humboldt State Geospatial Online. "Natural and False Color Composites." 2014. https://gsp.humboldt.edu/olm/Courses/GSP\_216/lessons/composites.html

\subsection{Độ phân giải không gian}

Khi một vệ tinh thu thập thông tin từ một khu vực quan sát, nó chia khu vực đó thành một ma trận hay một lưới. Mỗi ô trong ma trận được gọi là một \textbf{điểm ảnh (pixel)}. Mỗi ô là một hình vuông. Nếu ô vuông có kích thước 30 mét x 30 mét thì độ phân giải là 30 mét. Con số này biểu thị kích thước của một điểm ảnh. Hình 5 minh họa khái niệm này.

\textbf{Hình 5: Độ phân giải không gian}

Từ Hình 5, ta thấy nếu kích thước của một điểm ảnh nhỏ hơn, chẳng hạn 10 mét thay vì 30 mét, thì có thể đặt nhiều điểm ảnh hơn vào cùng một khu vực. Điều này có nghĩa là ảnh sẽ có độ phân giải cao hơn. Trong mỗi điểm ảnh, vệ tinh sẽ đo cường độ bức xạ điện từ (EMR) của dải sóng mà điểm ảnh đó nằm trong và gán một con số cho điểm ảnh. Nếu các ảnh số của một dải sóng từ vệ tinh có độ phân giải cao hơn, điều đó có nghĩa là ảnh có nhiều điểm ảnh hơn và cũng chứa nhiều dữ liệu hơn. Khi đọc một ảnh có độ phân giải cao hơn, ảnh trông tinh xảo hơn, các đường thẳng và đường cong mượt mà. Khi đọc một ảnh có độ phân giải thấp hơn, ảnh trông bị vỡ hạt (pixelated) hoặc giống như một bức tranh khảm. Hình 6 minh họa sự khác biệt này.

\textbf{Hình 6: Ảnh của cùng một vị trí với các độ phân giải khác nhau}

Từ Hình 6, ta thấy bức ảnh bên trái có độ phân giải 30 mét. Ta có thể quan sát nhiều chi tiết của khu vực từ bức ảnh này, như đường xá và các công trình. Tuy nhiên, bức ảnh bên phải có độ phân giải 300 mét. Bức ảnh trông giống một bức tranh khảm: ta chỉ thấy hình dạng mờ nhạt của khu vực và chỉ có thể phân biệt các khối màu của những đặc điểm địa lý như mặt nước và đất liền.

Hãy quay lại Hình 3 để xem cấu trúc độ phân giải của vệ tinh Landsat 8. Cột bên phải của hình cho thấy thiết lập độ phân giải của từng dải sóng trên Landsat 8. Độ phân giải của hầu hết các dải sóng là 30 mét. Tuy nhiên, Dải sóng 8 có độ phân giải 15 mét, trong khi Dải sóng 10 và Dải sóng 11 có độ phân giải 100 mét. Từ đó, ta thấy các dải sóng khác nhau trong cùng một vệ tinh có thể có độ phân giải khác nhau.

\subsection{Các vệ tinh phổ biến cung cấp ảnh}

Chuỗi Landsat của Hoa Kỳ và Sentinel 2 của Cơ quan Vũ trụ Châu Âu là hai hệ thống vệ tinh phổ biến nhất cung cấp ảnh vệ tinh phục vụ nghiên cứu. Vệ tinh mới nhất trong chuỗi Landsat là Landsat 9. Ngoài ra còn có các vệ tinh Landsat khác đang hoạt động, như Landsat 8 và Landsat 7. Sentinel 2 hiện có một cặp vệ tinh đang vận hành.

Các vệ tinh quay quanh Trái Đất để thu thập dữ liệu. Do chúng quay quanh Trái Đất với tốc độ khác với tốc độ tự quay của Trái Đất, nên chúng không thu thập dữ liệu từ cùng một vị trí mỗi ngày. Landsat 8 và Landsat 9 cùng nhau thu thập thông tin từ cùng một vị trí cứ sau 8 ngày. Cặp vệ tinh Sentinel 2 thu thập thông tin từ cùng một vị trí cứ sau 5 ngày. Số ngày mà một vệ tinh cần để thu thập thông tin lặp lại từ cùng một vị trí được gọi là \textbf{độ phân giải thời gian} (temporal resolution). Vui lòng xem video trong danh sách Tài liệu đọc bắt buộc để biết thêm thông tin về các vệ tinh thuộc chuỗi Landsat.

\subsection{Tóm tắt các kiến thức cơ bản về ảnh vệ tinh}

Hãy tóm tắt những gì chúng ta đã học cho đến nay. Chúng ta biết rằng vệ tinh ghi lại năng lượng phản xạ có bước sóng khác nhau từ bề mặt Trái Đất, dựa trên cấu trúc các băng (band) của nó. Với mỗi băng, dữ liệu thu thập được sẽ được lưu trong một điểm ảnh (pixel) của ma trận biểu diễn khu vực quan sát tại một thời điểm nhất định, và được tạo thành ảnh số của khu vực quan sát đó. Ảnh số của các băng khác nhau có thể được kết hợp để phục vụ nghiên cứu và phân tích. Do đó, điều quan trọng là phải biết băng nào hoặc tổ hợp băng nào cho ra ảnh phù hợp với nghiên cứu của bạn. Nếu nghiên cứu của bạn tập trung vào việc nhận dạng các đối tượng trên bề mặt Trái Đất, bạn cũng sẽ cần ảnh có độ phân giải cao cho nhiệm vụ này.

Giờ đây chúng ta đã có hiểu biết cơ bản về ảnh vệ tinh. Trong phần tiếp theo, chúng ta sẽ trình bày ngắn gọn về các xu hướng hiện nay trong việc sử dụng ảnh vệ tinh để thực hiện phân tích tài chính.

\section{Ứng dụng ảnh vệ tinh: Google Earth Engine}

Trong phần này, chúng ta sẽ sử dụng nền tảng Google Earth Engine (GEE) để truy xuất ảnh vệ tinh và thực hiện một số phân tích. Google Earth Engine là một nền tảng cho phép chúng ta phân tích ảnh vệ tinh và các dữ liệu không gian địa lý khác. Nhóm GEE đã thu thập và xử lý các ảnh vệ tinh lịch sử công khai từ các vệ tinh phổ biến, đồng thời lưu trữ chúng trong danh mục dữ liệu Earth Engine (Earth Engine data catalog). Do đó, chúng ta không cần tốn thời gian và công sức để thu thập và xử lý các ảnh vệ tinh lịch sử thô. GEE có giao diện trình soạn thảo mã (code editor) cho phép người dùng truy xuất và phân tích ảnh vệ tinh. GEE cũng có một API Python. Trong bài học này, chúng ta sẽ sử dụng API Python để truy xuất ảnh và thực hiện phân tích. Chúng ta sẽ dùng video của Văn phòng Điện toán Nghiên cứu UCLA (UCLA Office of Research Computing) (UCLA 2022) để minh họa ứng dụng này.

\subsection{Kịch bản: Vụ cháy Camp Fire năm 2018 ở Bắc California}

Vào thứ Năm, ngày 8 tháng 11 năm 2018, một đường dây truyền tải điện bị lỗi đã bốc cháy tại quận Butte, California, Hoa Kỳ. Do gió thổi xuống sườn dốc mạnh, đám cháy lan rất nhanh và thiêu rụi khoảng 153.000 mẫu Anh đất. Đám cháy chỉ dừng lại vào ngày 25 tháng 11 năm 2018. Camp Fire là thảm họa thiên nhiên gây thiệt hại tốn kém nhất năm 2018. Trong phần này, chúng ta sẽ xem xét các ảnh vệ tinh trước, trong và sau vụ cháy. Chúng ta cũng sẽ thực hiện một phân tích để khảo sát tình trạng thảm thực vật trước và sau vụ cháy.

\subsection{Thiết lập Google Earth Engine}

Trước khi sử dụng Google Earth Engine, vui lòng đọc tệp được liên kết bên dưới. Tệp này trình bày từng bước cách thiết lập tài khoản Google Earth Engine.

https://docs.google.com/document/d/1xZwsOHp49aCQFsPzvP0Kd8SJGmwMZA9Q/edit?usp=drive\_link\&ouid=103990088807346702419\&rtpof=true\&sd=true

\section{Minh họa Python API trên Google Colab}

\subsection{Tải các thư viện cần thiết cho phần minh họa này}

Chúng ta sẽ dùng Google Colab để truy cập Python API của GEE cho phần minh họa này.

\begin{lstlisting}[language=Python,escapeinside={(*@}{@*)}]
# pandas (*@để@*) (*@xử@*) (*@lý@*) (*@dữ@*) (*@liệu@*) (*@dạng@*) (*@bảng@*) (*@và@*) geopandas (*@để@*) (*@xử@*) (*@lý@*) (*@dữ@*) (*@liệu@*) (*@không@*) gian (*@địa@*) (*@lý@*)
import pandas as pd
import geopandas as gpd

# earth engine
import ee

# (*@Bản@*) (*@đồ@*) Google earth engine
import geemap

# cho (*@phép@*) (*@hiển@*) (*@thị@*) (*@ảnh@*) trong notebook
from IPython.display import Image
\end{lstlisting}

\subsection{Truy cập Google Earth Engine}

Chúng ta sẽ dùng đoạn mã sau để truy cập API của Google Earth Engine. Hãy thay \textquotesingle ee-si-learning-001\textquotesingle{} bằng ID dự án của riêng bạn. ID dự án được tạo khi bạn đăng ký Google Earth Engine. Trong quá trình xác thực, một hộp thoại sẽ xuất hiện và yêu cầu bạn chọn tài khoản Google đã dùng để đăng ký Google Earth Engine. Hãy nhấp vào tài khoản đó. Sau đó, sẽ có một số thông tin cùng với nút "continue" ở phía dưới bên phải của hộp thoại. Hãy nhấp vào nút "continue". Tiếp theo sẽ có một thông báo khác, và bạn cần đọc qua rồi nhấn nút "continue" để hoàn tất quá trình xác thực. Tổng cộng bạn sẽ phải nhấn hai nút continue.

\begin{lstlisting}[language=Python,escapeinside={(*@}{@*)}]
# (*@Bắt@*) (*@đầu@*) (*@quá@*) (*@trình@*) (*@xác@*) (*@thực@*) Google Earth Engine
ee.Authenticate()

# (*@Khởi@*) (*@tạo@*) Google Earth Engine (*@với@*) ID (*@dự@*) (*@án@*) (*@bạn@*) (*@đã@*) (*@thiết@*) (*@lập@*) (*@ở@*) (*@bước@*) (*@tạo@*) (*@tài@*) (*@khoản@*)
ee.Initialize(project='[REPLACE PROJECT ID]')
\end{lstlisting}

\subsection{Thiết lập các tham số lọc}

Hãy xác định vị trí quan tâm bằng tọa độ vĩ độ và kinh độ, cùng với thời điểm bắt đầu và kết thúc của giai đoạn cần lấy ảnh vệ tinh. Chúng ta có thể dùng Google Maps để tìm tọa độ vĩ độ và kinh độ của Quận Butte ở California. Vì đám cháy bắt đầu vào ngày 8 tháng 11 năm 2018, chúng ta sẽ lấy ảnh từ tháng 10 năm 2018 đến tháng 12 năm 2018.

\begin{lstlisting}[language=Python,escapeinside={(*@}{@*)}]
# (*@tọa@*) (*@độ@*) (*@của@*) (*@đám@*) (*@cháy@*) Camp Fire
lat =  39.444012
lon = -121.833619

# (*@điểm@*) quan (*@tâm@*) (*@dưới@*) (*@dạng@*) ee.Geometry
poi = ee.Geometry.Point(lon,lat)

# (*@ngày@*) (*@bắt@*) (*@đầu@*) (*@của@*) (*@khoảng@*) (*@thời@*) gian (*@cần@*) (*@lọc@*)
start_date = '2018-10-01'

# (*@ngày@*) (*@kết@*) (*@thúc@*)
end_date = '2019-01-31'
\end{lstlisting}

\subsection{Truy xuất ảnh từ Danh mục dữ liệu của Google Earth Engine}

Chúng ta sẽ truy xuất ảnh từ Landsat 8. Liên kết sau của Google Earth Engine cung cấp thông tin về bộ sưu tập ảnh này của Landsat 8. Trang này cũng cung cấp đoạn mã Python để lấy ảnh như được trình bày trong đoạn mã dưới đây.

https://developers.google.com/earth-engine/datasets/catalog/LANDSAT\_LC08\_C02\_T1\_L2

\begin{lstlisting}[language=Python,escapeinside={(*@}{@*)}]
# (*@lấy@*) (*@dữ@*) (*@liệu@*) (*@vệ@*) tinh
landsat = ee.ImageCollection("LANDSAT/LC08/C02/T1_L2")\
            .filterBounds(poi)\
            .filterDate(start_date,end_date)
\end{lstlisting}

\subsection{Kiểm tra thông tin ảnh}

Bây giờ chúng ta đã tải xong các ảnh. Hãy xem xét một số thông tin. Trước hết, hãy xem chúng ta nhận được bao nhiêu ảnh trong giai đoạn đã chọn. Độ phân giải thời gian của Landsat 8 là 16 ngày. Giai đoạn chúng ta chọn kéo dài 3 tháng. Do đó, chúng ta sẽ nhận được ít nhất 6 ảnh. ((3*30)/16 = 5,6)

\begin{lstlisting}[language=Python,escapeinside={(*@}{@*)}]
# (*@Kiểm@*) tra (*@số@*) (*@ảnh@*) (*@nhận@*) (*@được@*) trong giai (*@đoạn@*) (*@đã@*) (*@chọn@*).
print('Total number:', landsat.size().getInfo())
\end{lstlisting}

Tiếp theo, chúng ta dùng ảnh đầu tiên vừa lấy để xem có thể thu được những thông tin gì.

\begin{lstlisting}[language=Python,escapeinside={(*@}{@*)}]
landsat.first().getInfo()
\end{lstlisting}

Một yếu tố then chốt của ảnh vệ tinh là mức độ ảnh bị mây che phủ. Ảnh có độ che phủ mây dày sẽ không cung cấp nhiều thông tin cho phân tích của chúng ta. Thông thường, chúng ta muốn độ che phủ mây của ảnh vào khoảng hoặc dưới 0,05. Hãy xem thang đo độ che phủ mây của ảnh đầu tiên.

\begin{lstlisting}[language=Python,escapeinside={(*@}{@*)}]
landsat.first().get('CLOUD_COVER').getInfo()
\end{lstlisting}

Chúng ta thấy độ che phủ mây của ảnh đầu tiên là 0,05, đây là mức tốt. Bây giờ hãy kiểm tra xem ảnh đầu tiên được chụp khi nào.

\begin{lstlisting}[language=Python,escapeinside={(*@}{@*)}]
landsat.first().get('DATE_ACQUIRED').getInfo()
\end{lstlisting}

Điều tiếp theo chúng ta có thể kiểm tra là tên các dải phổ (band) thu được.

\begin{lstlisting}[language=Python,escapeinside={(*@}{@*)}]
landsat.first().bandNames().getInfo()
\end{lstlisting}

Từ kết quả của đoạn mã trên, chúng ta thấy có nhiều hơn 11 dải phổ so với những gì đã học ở phần trước. Các tên dải phổ có \textquotesingle ST\_\textquotesingle{} đều thuộc về một dải phổ. Chúng là các dải phổ con của dải Aerosol.

\subsection{Trực quan hóa ảnh}

Giờ đây chúng ta đã sẵn sàng xem các ảnh vừa lấy. Trước hết, hãy tạo nhãn cho các ảnh để dùng trong phần trực quan hóa tiếp theo.

\begin{lstlisting}[language=Python,escapeinside={(*@}{@*)}]
# (*@đưa@*) (*@các@*) (*@ảnh@*) (*@vào@*) (*@một@*) danh (*@sách@*)
landsat_list = landsat.toList(landsat.size());

#(*@Tạo@*) (*@nhãn@*) cho (*@các@*) (*@ảnh@*)
labels = ["Image #"+str(i)+", "+str(ee.Image(landsat_list.get(i)).get('DATE_ACQUIRED').getInfo())+", Cloud cover:"+ str(ee.Image(landsat_list.get(i)).get('CLOUD_COVER').getInfo()) for i in range(landsat.size().getInfo())]
labels
\end{lstlisting}

Từ các nhãn ảnh, chúng ta thấy có 8 ảnh, từ ảnh 0 đến ảnh 7. Chúng ta cũng thấy ngày chụp và chỉ số che phủ mây của từng ảnh.

Bây giờ hãy xác định một số tham số để hiển thị các ảnh. Chúng ta sẽ dùng các dải đỏ, lục và lam để tạo thành một ảnh giống ảnh chụp thông thường. Chúng ta cũng xác định kích thước điểm ảnh (pixel) sẽ hiển thị. Cuối cùng, chúng ta đặt độ sáng bằng cách gán các giá trị min và max. Việc gán giá trị cho min và max là một quá trình thử và sai. Bạn cần thử nghiệm các con số để có được độ sáng phù hợp cho ảnh.

\begin{lstlisting}[language=Python,escapeinside={(*@}{@*)}]
parameters = {
                'min': 7000,
                'max': 16000,
                'dimensions': 800, # (*@kích@*) (*@thước@*) (*@cạnh@*) (*@hình@*) (*@vuông@*) (*@tính@*) (*@bằng@*) pixel
                'bands': ['SR_B4', 'SR_B3', 'SR_B2'] # (*@các@*) (*@dải@*) (*@phổ@*) (*@để@*) (*@hiển@*) (*@thị@*) (r,g,b)
             }
\end{lstlisting}

Bây giờ hãy dùng geemap để hiển thị ảnh đầu tiên chồng lên bản đồ.

\begin{lstlisting}[language=Python,escapeinside={(*@}{@*)}]
Map = geemap.Map()

first_image = landsat.first()

Map.addLayer(first_image, parameters, "First_Image")
Map.setCenter(lon,lat,8)
Map
\end{lstlisting}

Từ bản đồ trên, bạn có thể thấy ảnh đầu tiên trong bộ sưu tập ảnh của chúng ta. Đây là ảnh chụp ngày 7 tháng 11 năm 2018. Bạn có thể di chuột lên bản đồ và kéo bản đồ đi các hướng. Bạn có thể phóng to hoặc thu nhỏ ảnh bằng cách nhấp vào biểu tượng "+" hoặc "-" ở phía bên trái của bản đồ.

Ta vừa thấy cách hiển thị một ảnh trên bản đồ. Để hiển thị tương tác các ảnh khác trong bộ sưu tập, ta tạo một bản đồ mới và thêm thanh trượt chuỗi thời gian (time series slider) lên bản đồ.

\begin{lstlisting}[language=Python,escapeinside={(*@}{@*)}]
image = landsat.toBands()

Map2 = geemap.Map()
Map2.addLayer(image,{},"Time series", False)
Map2.setCenter(lon,lat,8)
Map2.add_time_slider(landsat, parameters, labels = labels, time_interval=1)
Map2
\end{lstlisting}

Bản đồ mới có thanh trượt ở góc dưới bên phải. \textbf{Lưu ý khi dùng thanh trượt:} do một số vấn đề kỹ thuật, không dùng nút "Play the time slider".

Cạnh thanh trượt hiển thị nhãn của ảnh đầu tiên (ảnh \#0). Di chuột tới chấm bên trái thanh trượt, chấm sẽ chuyển xanh; kéo chấm xanh dọc thanh để xem các ảnh khác. Nhãn cho biết thông tin ảnh đang hiển thị. Chẳng hạn, ảnh \#3 có nhiều mây (độ phủ mây 67) nên phần lớn là màu trắng; ảnh \#5 có độ phủ mây 6 nên thấy rõ bề mặt Trái Đất hơn.

\subsection{Chọn và phóng to các ảnh không mây}

Sau khi xem toàn bộ ảnh lấy từ danh mục dữ liệu, ta chọn ảnh dùng cho phân tích: 3 ảnh không mây, gồm một ảnh trước đám cháy, một ảnh trong đám cháy và một ảnh sau đám cháy, đồng thời phóng to vùng quan tâm. Qua xem xét 8 ảnh, ảnh \#0, \#2 và \#5 đáp ứng tiêu chí.

\begin{lstlisting}[language=Python,escapeinside={(*@}{@*)}]
#Select images we want and create a new image collection
img1 = ee.Image(landsat_list.get(0))
img2 = ee.Image(landsat_list.get(2))
img3 = ee.Image(landsat_list.get(5))

landsat_2 = ee.ImageCollection.fromImages([img1, img2, img3])
print('Total number:', landsat_2.size().getInfo())
\end{lstlisting}

Tiếp theo, tạo nhãn mới cho bộ ảnh mới và cập nhật tham số hiển thị.

\begin{lstlisting}[language=Python,escapeinside={(*@}{@*)}]
# Create a new list
landsat_list_2 = landsat_2.toList(landsat_2.size())

# Define a region of interest with a buffer zone
roi = poi.buffer(20000) # meters

# New visualization parameters
parameters_2 = {
                'min': 6000,
                'max': 16000,
                'dimensions': 800,
                'bands': ['SR_B4', 'SR_B3', 'SR_B2'],
                'region':roi
             }

#Create new labels for images
labels_2 = ["Image #"+str(i)+", "+str(ee.Image(landsat_list_2.get(i)).get('DATE_ACQUIRED').getInfo())+", Cloud cover:"+ str(ee.Image(landsat_list_2.get(i)).get('CLOUD_COVER').getInfo()) for i in range(landsat_2.size().getInfo())]
labels_2
\end{lstlisting}

\begin{lstlisting}[language=Python,escapeinside={(*@}{@*)}]
image2 = landsat_2.toBands()

Map3 = geemap.Map()
Map3.addLayer(image2,{},"Time series", False)
Map3.setCenter(lon,lat,8)
Map3.add_time_slider(landsat_2, parameters_2, labels = labels_2, time_interval=1)
Map3
\end{lstlisting}

Góc trên của ảnh giữa có thấy đám cháy, nhưng so sánh ảnh đầu với ảnh cuối vẫn khó nhận ra tác động của đám cháy. Ta dùng chỉ số NDVI để đánh giá thiệt hại rõ hơn.

\subsection{Chỉ số thực vật khác biệt chuẩn hóa (NDVI)}

Chỉ số thực vật khác biệt chuẩn hóa (Normalized Difference Vegetation Index, NDVI) dùng để đánh giá sức khỏe thực vật, được tính từ ánh sáng phản xạ ở băng đỏ và băng cận hồng ngoại của ảnh vệ tinh. Chỉ số càng cao thì cây càng khỏe và xanh; chỉ số càng thấp thì cây càng úa và kém khỏe. Khái niệm được minh họa ở Hình 7.

\textbf{Hình 7. Minh họa NDVI}

Ở phần sau, ta chuyển ba ảnh thành ảnh NDVI. Mỗi pixel có một giá trị NDVI: giá trị cao được tô xanh lá, giá trị thấp tô đỏ, giá trị trung gian tô vàng.

Trước hết, tính giá trị NDVI cho từng ảnh.

\begin{lstlisting}[language=Python,escapeinside={(*@}{@*)}]
#Calculate NDVI for each image
img_ndvi_1 = ee.Image(landsat_list_2.get(0)).normalizedDifference(['SR_B5', 'SR_B4'])
img_ndvi_2 = ee.Image(landsat_list_2.get(1)).normalizedDifference(['SR_B5', 'SR_B4'])
img_ndvi_3 = ee.Image(landsat_list_2.get(2)).normalizedDifference(['SR_B5', 'SR_B4'])

landsat_3 = ee.ImageCollection.fromImages([img_ndvi_1, img_ndvi_2, img_ndvi_3])
print('Total number:', landsat_3.size().getInfo())
\end{lstlisting}

Sau đó, thiết lập bảng màu mới và tham số hiển thị NDVI.

\begin{lstlisting}[language=Python,escapeinside={(*@}{@*)}]
# Create a new list
landsat_list_3 = landsat_3.toList(landsat_3.size())

# ndvi palette: red is low, green is high vegetation
palette = ['red', 'yellow', 'green']

ndvi_parameters = {'min': 0,
                   'max': 0.4,
                   'dimensions': 512,
                   'palette': palette,
                   'region': roi}
\end{lstlisting}

Bây giờ hiển thị các ảnh NDVI.

\begin{lstlisting}[language=Python,escapeinside={(*@}{@*)}]
image3 = landsat_3.toBands()

Map4 = geemap.Map()
Map4.addLayer(image3,{},"Time series", False)
Map4.setCenter(lon,lat,8)
Map4.add_time_slider(landsat_3, ndvi_parameters, labels = labels_2, time_interval=1)
Map4
\end{lstlisting}

Xét phần giữa phía trên của mỗi ảnh NDVI tương tác. Ảnh \#0 (trước Camp Fire) có xen lẫn xanh, vàng và đỏ. Ảnh \#1 (trong Camp Fire) có một dải đỏ lớn chạy từ phải sang trái rồi đi xuống; vùng đỏ là nơi lửa cháy. Ảnh \#3 (sau Camp Fire) có vùng đỏ lớn ở phần giữa phía trên và một phần ở phía trên bên trái. Tóm lại, ảnh NDVI giúp nhận diện rõ hơn các thay đổi của thảm thực vật trên bề mặt Trái Đất.

\section{Kết luận}

Trong bài học này, trước hết chúng ta tìm hiểu kiến thức cơ bản về viễn thám, phương pháp thu nhận ảnh vệ tinh. Tiếp đó, chúng ta đi qua các nội dung nền tảng về ảnh vệ tinh, gồm phổ điện từ, độ phân giải không gian và các vệ tinh thường dùng. Sau đó, bài học giới thiệu một số trường hợp ứng dụng ảnh vệ tinh trong phân tích tài chính. Cuối cùng, chúng ta dùng Python API của Google Earth Engine để tải một số ảnh vệ tinh và thực hiện phân tích; cụ thể là dùng ảnh NDVI để khảo sát mức độ thiệt hại thảm thực vật sau vụ cháy Camp Fire năm 2018 ở California.
